\subsection*{Independent second reading of the canonical-point criterion}

All 223 lines of \texttt{finite\_checks/b1\_canonical\_point.tex}
were independently read on 18 November 2026.  The outcome for
B.141 is \textbf{PASS}, with no correction requested.  This is
a separate same-model-family check by GPT-6 Astra Pro in Pro
mode; a separate effort setting is not exposed.  It does not
assert a human check or priority.

The constant Lang quotient is pointed and geometrically
connected for the reason stated: the nonzero prime-to-$p$
abelian quotient remains nonzero after Abel--Jacobi pullback.
For genus two, Riemann--Roch gives one effective divisor in
each noncanonical rational degree-two class and $q+1$ in
the canonical class.  Thus subtracting the common character
sum gives exactly $q\zeta^{i\chi(K_C-2P)}$ as the last
coefficient.  For larger degrees the common section count
makes every coefficient vanish.  The rational marking
justifies the use of actual rational line-bundle classes.

The twist convention is consistent at every closed point:
degree $m$ adds the translation $mt$, so the linear and
quadratic coefficients acquire respectively $\zeta^{it}$
and $\zeta^{2it}$.  The character factorization gives the
displayed rational integral class-number ratio.

The canonical degree-two map is separable also in
characteristic two.  Over a perfect field a purely
inseparable quadratic extension of a rational function
field is again rational, which would contradict genus two.
The involution therefore exists over $k$ and gives
$S_{-i}=\zeta^{-ia}S_i$.  Even order of $p$ modulo the odd
prime $\ell$ puts $-1$ in the subgroup generated by $p$.
Reduction of a vanishing character factor and its
$p^{f/2}$-power then gives precisely $a+2t=0$.
This pairing works simultaneously for every primitive
character and also when $p=2$; the sign reduces harmlessly.

The count of covers divides the number of nonzero
functionals on the marked vector by $\ell-1$.  Distinct
resulting kernels give distinct pointed covers: any
arithmetic character with trivial geometric restriction
and split marked fibre would itself be trivial.  Hence
the geometric quotient is not being overcounted by
different constant twists.

For constant extension of degree prime to $\ell$, the
base-changed character is the composition with the
Jacobian norm.  This follows directly from the factorization
of $F_q^n-1$.  The norm multiplies the rational canonical
class by $n$, so nonvanishing persists for the very same
cover, not just for an independently chosen new cover.
Conversely a divisibility relation already over $k$
remains one over every constant extension.

The statement is a sufficient condition on the specified
marking.  It neither replaces nonvanishing modulo
$\ell J(k)$ by mere nonzero primary torsion nor excludes
the zero canonical-point class.  Its remaining-locus
paragraph is therefore logically accurate.  Nonemptiness
of this sufficient class is a separate example calculation
and was not presumed in the first reading.

\medskip\noindent\textbf{Separate check of the explicit-family addition.}
The subsequent addition of $C_p:y^2=x^6+1$ with $P=(0,1)$,
for primes $p\ge5$ congruent to $2$ or $5$ modulo nine,
was then fully read and also \textbf{passed}.  The map
$(x,y)\mapsto(x^2,y)$ is a degree-two morphism to
$E_p:v^2=u^3+1$, sends both infinities to the origin,
and sends the canonical-point class to $-2R=R$ for the
nonzero flex $R=(0,1)$ of order three.  The cubing
permutation of $\mathbb F_p$ gives $\#E_p=p+1$, whose
three-adic valuation is one.  Thus $R$ is not divisible
by three, and the claimed nondivisibility of the
canonical-point class follows by pushforward.

The complete finite certificate code was read, including
its finite-field operations, elliptic addition, character
table, and both polynomial calculations.  As an
independent numerical check, $\mathbb F_{25}^{\times}$
is cyclic of order $24$, and $x\mapsto x^6$ has image
$\mu_4=\mathbb F_5^{\times}$ with six preimages each.
The values of $1+\mu_4$ are $2,3,4,0$, with all three
nonzero values squares in $\mathbb F_{25}$.  Hence
the affine point count is $2+6(2+2+2+1)=44$, giving
$46$ after the two points at infinity.  The six
$\mathbb F_5$-points have the asserted character
occupancies $[1,1,4]$.  Multiplying the two conjugate
quadratics gives $1-3T+4T^2-15T^3+25T^4$ and value
$12$ at one, agreeing with the Jacobian orders
$36$ and $432$.  This additional event is distinct
from the first full reading of the conditional criterion.
